\ulohahead{společná informatika \textnormal{\footnotesize[úroveň 2]}}{Pumping lemma pro bezkontextové jazyky a Master theorem}
\begin{zadani}
\begin{enumerate}[label=(\alph*)]
  \item \bod{1} Zformulujte pumping lemma pro bezkontextové jazyky (CFL).
  \item \bod{2} Dokažte, že $L=\{a^n b^n c^n\mid n\ge0\}$ není bezkontextový.
  \item \bod{3} Vysvětlete hlavní myšlenku Master theoremu (rekurzní strom) a jeho tři případy. \emph{(Požadavek uvádí Master theorem ,,bez důkazu'' - formální důkaz se nevyžaduje, jen pochopení.)}
\end{enumerate}
\end{zadani}

\begin{reseni}
\textbf{(a)} Pro každý CFL existuje $p$ tak, že každé slovo $w\in L$ s $|w|\ge p$ lze psát $w=uvxyz$, kde $|vxy|\le p$, $|vy|\ge1$ a pro každé $i\ge0$ je $uv^ixy^iz\in L$. (Pumpují se \emph{dva} úseky najednou.)

\textbf{(b)} Sporem: nechť $L$ je CFL s konstantou $p$ a $w=a^pb^pc^p$. Rozklad $w=uvxyz$ s $|vxy|\le p$ znamená, že $vxy$ zasahuje nejvýše \emph{dva} ze tří bloku ($a,b,c$) - nemůže obsahovat $a$ i $c$, neboť mezi nimi je $p$ symbolu $b$. Napumpováním ($i=2$) se zvýší počet symbolu jen v jednom či dvou blocích, takže přestanou být všechny tři stejně početné: $uv^2xy^2z\notin L$. Spor, $L$ není CFL.

\textbf{(c)} \emph{Master theorem} ($T(n)=aT(n/b)+f(n)$) přes \emph{rekurzní strom}: na úrovni $i$ je $a^i$ podproblému velikosti $n/b^i$, práce mimo rekurzi na úrovni $i$ je $a^i f(n/b^i)$; listu je $a^{\log_b n}=n^{\log_b a}$. Sečtením úrovní rozhoduje, zda dominují listy ($n^{\log_b a}$), kořen ($f(n)$), nebo jsou úrovně vyrovnané (pak faktor $\log n$). Konkrétně pro $f(n)=\Theta(n^c)$ ($a\ge1,b>1$) porovnej $c$ s $\log_b a$: $c<\log_b a\Rightarrow\Theta(n^{\log_b a})$ (listy); $c=\log_b a\Rightarrow\Theta(n^{c}\log n)$; $c>\log_b a\Rightarrow\Theta(n^{c})$ (kořen, za podmínky regularity). To jsou tři případy věty.
\end{reseni}

\begin{jadro}
CFL pumping: $w=uvxyz$, $|vxy|\le p$, $|vy|\ge1$, $uv^ixy^iz\in L$. $\{a^nb^nc^n\}$: $vxy$ pokryje max 2 bloky $\Rightarrow$ pumpování rozváží třetí. Master theorem: porovnej $n^{\log_b a}$ (listy) s $f(n)$.
\end{jadro}

\kalibrace{Důkaz nebezkontextovosti (přes pumping) podpírá zařazení do Chomského hierarchie - to je v požadavku. Master theorem se dle požadavku umí \emph{použít} (bez důkazu); intuice rekurzního stromu je jen pro pochopení.}
\past{CFL pumping pumpuje \emph{dva} úseky ($v$ a $y$) zároveň, na rozdíl od regulárního (jeden). Klíč u $a^nb^nc^n$ je $|vxy|\le p$ (nepokryje krajní bloky současně).}
